% Introduction

日本のアニメ産業は、かつてない商業的成功を収めている。
国内アニメ制作市場は2024年に過去最高の3{,}621億円に達し~\cite{tdb2025anime}、
世界のアニメ市場は年平均成長率10.6\%で成長し、
2031年までに496億ドルに達すると予測されている~\cite{mordor2026anime}。
しかし、この成長の背後には根深い構造的問題が存在する。
アニメを制作するクリエイティブ人材は、依然としてその大部分が不可視のままである。

アニメ制作労働の極めて大きな割合をフリーランスが担っている。
日本アニメーター・演出協会（JAniCA）によれば、
2023年時点でアニメーション制作者の47.3\%が個人事業主であり~\cite{janica2023survey}、
これは全国平均の約6倍にあたる。
これらのフリーランスは、スタジオが個々のエピソードや特定の制作タスクを
二次・三次の下請けに日常的に外注する多層的な下請けシステムの中で活動している。
このシステムにおいて、個々のアニメーターはしばしばクレジットされない。
最終クレジットには下請けスタジオの名前のみが記載されるか、
あるいは一切の帰属表示がなされない場合もある。

この\textit{貢献の不可視性問題}（contribution invisibility problem）は具体的な弊害をもたらしている。
フリーランスアニメーターは、将来のクライアントに対して自身の作業履歴を
検証可能な形で示すことができず、
労働者の37.7\%が月収20万円未満である~\cite{nafca2024labor}業界において
交渉力が弱体化している。
スタジオが閉鎖した場合---これは頻繁に起こる---、
誰がどの作品に携わったかという非公式な記録は完全に消失する。
日本のフリーランス保護法（2024年）や、
公正取引委員会によるアニメ制作業界の労働慣行に関する調査~\cite{jftc2025anime}
といった近年の政策的対応は、契約上の保護を扱ってはいるものの、
帰属表示の問題を解決するものではない。

これと並行して、ブロックチェーン技術は分散型資格証明のための
新たなプリミティブ（primitive）を生み出してきた。
Weyl et al.~\cite{buterin2022decentralized}は、
\textbf{ソウルバウンドトークン（Soulbound Token, SBT）}---
特定のブロックチェーンアドレスに紐づけられた譲渡不可能なトークン---を、
従来のNFTが持つ投機的性質を伴うことなく、
資格、コミットメント、および所属をエンコードする仕組みとして提案した。
SBTはその後、医療資格証明、学術証明書の検証、農業サプライチェーンなどに適用されてきたが、
\textbf{クリエイティブ制作ワークフローやフリーランスアニメーターの資格証明に
SBTを適用した先行研究は存在しない。}

注目すべきは、アニメ産業における既存のブロックチェーン応用が、
ほぼ例外なく\textit{ファン向け}のマネタイズに焦点を当ててきたことである。
IPのトークン化、NFTコレクティブル、そしてオーディエンスから価値を引き出すための
ファンエンゲージメントプラットフォームがその例である~\cite{fourpillars2025anime,animechain2024whitepaper}。
これに対し、本論文ではブロックチェーンを\textit{労働者保護}のツールとして用いる。
ファンから収益を得るためではなく、個々のクリエイターの
職業的アイデンティティとキャリア記録を保護するためにブロックチェーンを活用する。

本論文は、フリーランスアニメーター向けの既存の制作管理ツールに
SBT発行を統合するよう設計されたブロックチェーンベースの貢献証明フレームワーク
\textbf{AnimatorSBT}を提示する。
本システムは、個々の制作タスク
（例：原画カット、動画フレーム、仕上げ作業）の粒度で、
不変かつ検証可能で可搬性のある貢献記録を作成する。
重要な点として、本フレームワークは\textit{エンドユーザーにとって透過的}であるよう設計されている。
アニメーターは管理ツールにタスクを記録するという通常のワークフローを継続し、
プロジェクトのマイルストーン到達時に貢献証明書が自動的に発行される。
ブロックチェーンの知識、ウォレット管理、ガス代の支払いは一切不要である。

本論文の貢献は以下の通りである。

\begin{enumerate}
  \item アニメ産業における\textbf{貢献の不可視性問題}を形式化する。
    これを実際の貢献者の集合とクレジットされた貢献者の集合との間のギャップとして定義し、
    理想的なソリューションのための5つの要件を確立する
    （\cref{sec:problem}）。
  \item SBTスマートコントラクト、発行サービス、
    および制作管理ツールと統合された検証ポータルから構成される
    \textbf{AnimatorSBTフレームワーク}を提案する
    （\cref{sec:framework}）。
  \item スマートコントラクト、発行サービス、および検証ポータルを含む、
    \textbf{Polygon Amoyテストネット上にデプロイされた動作するプロトタイプ}を提示する。
    実測されたガスコストを用いたコストシミュレーションにより、
    産業全体での採用コストは年間約\textbf{\$256}---
    産業の年間売上高の$0.000011\%$---であることを示す（\cref{sec:evaluation}）。
  \item スタジオの協力に関する課題、自己申告のリスク、
    および日本固有の規制上の考慮事項を含む\textbf{採用障壁と悪用ベクトル}を特定し、
    緩和策を提案する（\cref{sec:discussion}）。
\end{enumerate}

本論文の構成は以下の通りである。
\Cref{sec:related_work}では、アニメ制作、労働環境、SBT、
およびクリエイティブ経済におけるブロックチェーンに関する関連研究を調査する。
\Cref{sec:problem}では、問題を形式化し、既存のソリューションを評価する。
\Cref{sec:framework}では、AnimatorSBTフレームワークを記述する。
\Cref{sec:implementation}では、技術仕様とプロトタイプの実装を提示する。
\Cref{sec:evaluation}では、コストシミュレーションと要件評価を提示する。
\Cref{sec:discussion}では、実践的な課題と一般化可能性について議論する。
\Cref{sec:conclusion}では、今後の研究の方向性とともに結論を述べる。
